\documentclass[tikz, border=10mm]{standalone}
\usepackage[utf8]{inputenc}
\usepackage[T1]{fontenc}
\usepackage[french]{babel}
\usepackage{fontawesome5}
\usepackage{xcolor}
\usepackage{amssymb}
\usepackage{amsmath}

\renewcommand{\familydefault}{\sfdefault}
\usetikzlibrary{positioning, shadows, backgrounds, fit, babel}

\definecolor{myblue}{RGB}{41, 128, 185}
\definecolor{myteal}{RGB}{22, 160, 133}
\definecolor{mygreen}{RGB}{39, 174, 96}
\definecolor{mypurple}{RGB}{142, 68, 173}
\definecolor{myorange}{RGB}{211, 84, 0}
\definecolor{myred}{RGB}{210, 40, 40}

\tikzset{
    partbox/.style={fill=white, draw=#1, line width=1.5pt, rounded corners=6pt, inner sep=12pt, align=left, text width=13cm, drop shadow={opacity=0.08, shadow xshift=1.5mm, shadow yshift=-1.5mm}},
    chapbox/.style={fill=white, draw=gray!30, thick, rounded corners=4pt, inner sep=10pt, align=left, text width=12.5cm, drop shadow={opacity=0.04, shadow xshift=1.5mm, shadow yshift=-1.5mm}},
    milebox/.style={fill=myorange!5, draw=myorange, line width=2pt, rounded corners=6pt, inner sep=12pt, align=left, text width=13cm, drop shadow={opacity=0.08, shadow xshift=1.5mm, shadow yshift=-1.5mm}},
    partdate/.style={text=white, font=\bfseries\large, rounded corners=4pt, inner sep=8pt, align=center, minimum width=3.5cm},
    chapdate/.style={text=myred, font=\bfseries\small, align=right, minimum width=3cm},
    miledate/.style={fill=myorange, text=white, font=\bfseries\Large, rounded corners=4pt, inner sep=10pt, align=center, minimum width=3.5cm},
}

\newcommand{\bulletitem}[2][myblue]{%
    \par\hangindent=1.2em\noindent\makebox[1.2em][l]{\color{#1}$\bullet$}#2\par
}

\newcommand{\adddatepart}[3]{
    \node[partdate, fill=#3, anchor=east] (D#1) at (-0.5, 0 |- #1.center) {#2};
    \draw[draw=#3, line width=1.5pt] (D#1) -- (#1.west |- #1.center);
    \fill[#3] (0, 0 |- #1.center) circle (5pt);
}
\newcommand{\adddatechap}[2]{
    \node[chapdate, anchor=east] (D#1) at (-0.5, 0 |- #1.center) {#2};
    \draw[draw=myred!40, thick] (D#1) -- (#1.west |- #1.center);
    \fill[myred!80] (0, 0 |- #1.center) circle (3.5pt);
}
\newcommand{\adddatemile}[2]{
    \node[miledate, anchor=east] (D#1) at (-0.5, 0 |- #1.center) {#2};
    \draw[draw=myorange, line width=2pt] (D#1) -- (#1.west |- #1.center);
    \fill[myorange] (0, 0 |- #1.center) ++(0,-6pt) -- ++(6pt,6pt) -- ++(-6pt,6pt) -- ++(-6pt,-6pt) -- cycle;
}

\newcommand{\headerblock}{
\node[anchor=south west] at (-4.5, 1.5) {
    \begin{minipage}{15cm}
        \Huge\bfseries\color{myblue} Rétro-planning de Thèse \\[3mm]
        \Large\color{gray!80!black} Objectif final : \textbf{Soutenance le 8 septembre 2026}\\[2mm]
        \normalsize\color{myred}\faCalendarCheck\ \textit{Les dates en rouge correspondent aux livraisons prévues aux encadrants.}
    \end{minipage}
};
}

\begin{document}

% ================= PAGE 1 : PARTIE I =================
\begin{tikzpicture}
\headerblock

\node[partbox=myblue, anchor=north west] (P1) at (1, 0) {
    \textbf{\Large \color{myblue}Partie I : Du Single-Linkage au clustering hiérarchique}\\
    \vspace{1.5mm}
    {\small\color{myorange}\faInfoCircle\ \textit{Note : Partir du graphe général pour se restreindre ensuite au graphe complet des distances euclidiennes ($\mathbb{R}^d$).}}
};
\adddatepart{P1}{15 mars\\-- 10 avr.}{myblue}

\node[chapbox, below=0.5cm of P1.south west, anchor=north west] (C1) {
    \textbf{\large Chapitre 1 : Le Single-Linkage : trois points de vue équivalents}\\
    \vspace{1.5mm}
    {\normalsize\color{black!80}
    \bulletitem{Introduction au Single-Linkage}
    \bulletitem{Le Single-Linkage en pratique : l'arbre minimum couvrant}
    \bulletitem{Le Single-Linkage en théorie : persistance de graphes géométriques}
    \bulletitem{Le modèle statistique du Single-Linkage : Hartigan et les amas de forte densité}
    }
    \vspace{1.5mm}
    {\small\color{myorange}\faEdit\ \textit{Action : Refaire quelques figures. Pas de nouvelles expériences requises.}}
};
\adddatechap{C1}{15 -- 20 mars}

\node[chapbox, below=0.3cm of C1.south west, anchor=north west] (C2) {
    \textbf{\large Chapitre 2 : Les limites du Single-Linkage}\\
    \vspace{1.5mm}
    {\normalsize\color{black!80}
    \bulletitem{Quand la notion de densité manque : le \emph{chaining effect}}
    \bulletitem{L'état de l'art aujourd'hui : HDBSCAN}
    }
};
\adddatechap{C2}{20 -- 24 mars}

\node[chapbox, below=0.3cm of C2.south west, anchor=north west] (C3) {
    \textbf{\large Chapitre 3 : Une vue générale sur le clustering hiérarchique. Introduire des a priori géométriques. Ou comment « hacker » HDBSCAN}\\
    \vspace{1.5mm}
    {\small\color{myorange}\faEdit\ \textit{Action : Refaire des figures et mettre à jour les résultats sur SemanticKITTI.}}
};
\adddatechap{C3}{24 -- 30 mars}

\node[chapbox, below=0.3cm of C3.south west, anchor=north west] (C4) {
    \textbf{\large Chapitre 4 : Le phénomène de percolation au cœur des performances}\\
    \vspace{1.5mm}
    {\normalsize\color{black!80}
    \bulletitem{Introduction à la percolation}
    \bulletitem{La \emph{vitesse de percolation}}
    }
    \vspace{1.5mm}
    {\small\color{myorange}\faEdit\ \textit{Action : Figures illustratives. Calculer plus précisément la vitesse de percolation (dim 2,3, grands $K$).}}\\[0.5mm]
    {\small\color{myorange}\faQuestionCircle\ \textit{Question : Pertinent de placer ces résultats avant la $K$-généralisation du Single-Linkage ?}}
};
\adddatechap{C4}{30 mars\\-- 10 avr.}

\begin{scope}[on background layer]
    \draw[gray!20, line width=6pt] (0, 1.5) -- (0, 0 |- C4.center);
\end{scope}
\end{tikzpicture}

\newpage

% ================= PAGE 2 : PARTIE II =================
\begin{tikzpicture}
\headerblock

\node[partbox=myteal, anchor=north west] (P2) at (1, 0) {
    \textbf{\Large \color{myteal}Partie II : La généralisation naturelle au moyen d'hypergraphes}\\
    \vspace{1.5mm}
    {\small\color{myorange}\faEdit\ \textit{Action : Refaire quelques figures et étoffer la bibliographie (un peu de travail).}}
};
\adddatepart{P2}{10 avr.\\-- 1\textsuperscript{er} mai}{myteal}

\node[chapbox, below=0.5cm of P2.south west, anchor=north west] (C5) {
    \textbf{\large Chapitre 5 : Introduction : retour sur la non prise en compte de la densité du Single-Linkage}
};
\adddatechap{C5}{10 -- 15 avr.}

\node[chapbox, below=0.3cm of C5.south west, anchor=north west] (C6) {
    \textbf{\large Chapitre 6 : Hypergraph-Percol : la généralisation du Single-Linkage, compatible aux trois points de vue}\\
    \vspace{1.5mm}
    {\normalsize\color{black!80}
    \bulletitem[myteal]{La généralisation naturelle du graphe géométrique et de la notion de composante connexe : le complexe de \v{C}ech et les $K$-polyèdres}
    \bulletitem[myteal]{La correspondance avec le modèle de Hartigan et l'estimateur de densité des $K$-Plus Proches Voisins}
    }
};
\adddatechap{C6}{15 -- 20 avr.}

\node[chapbox, below=0.3cm of C6.south west, anchor=north west] (C7) {
    \textbf{\large Chapitre 7 : À la recherche de la généralisation de l'arbre minimum couvrant}\\
    \vspace{1.5mm}
    {\normalsize\color{black!80}
    \bulletitem[myteal]{L'approche persistante : les simplexes $K$-séparants}
    \bulletitem[myteal]{Une condition nécessaire : être $K$-Gabriel}
    \bulletitem[myteal]{La bonne structure à observer : la mosaïque de Delaunay d'ordre $K$}
    }
    \vspace{1.5mm}
    {\small\color{myorange}\faEdit\ \textit{Action : Ajouter quelques lemmes techniques (peu de travail).}}
};
\adddatechap{C7}{20 -- 25 avr.}

\node[chapbox, below=0.3cm of C7.south west, anchor=north west] (C8) {
    \textbf{\large Chapitre 8 : Conclusion sur la partie : la richesse structurelle de la mosaïque d'ordre $K$}\\
    \vspace{1.5mm}
    {\small\color{myorange}\faInfoCircle\ \textit{Note : Ouverture sur d'autres questions. Reprendre le rapport de recherche (rester vague).}}
};
\adddatechap{C8}{25 avr.\\-- 1\textsuperscript{er} mai}

\begin{scope}[on background layer]
    \draw[gray!20, line width=6pt] (0, 1.5) -- (0, 0 |- C8.center);
\end{scope}
\end{tikzpicture}

\newpage

% ================= PAGE 3 : PARTIE III =================
\begin{tikzpicture}
\headerblock

\node[partbox=mygreen, anchor=north west] (P3) at (1, 0) {
    \textbf{\Large \color{mygreen}Partie III : Vers des données davantage structurées}\\
    \vspace{1.5mm}
    {\small\color{myorange}\faEdit\ \textit{Action : Faire 2 sous-parties distinctes (2 intros ?). Reprendre les articles (Nahuel \& Cerema) $\rightarrow$ peu de travail s'ils sont aboutis.}}
};
\adddatepart{P3}{1\textsuperscript{er} -- 20 mai}{mygreen}

\node[chapbox, below=0.5cm of P3.south west, anchor=north west] (C9) {
    \textbf{\large Chapitre 9 : Introduction}\\
    \vspace{1.5mm}
    {\normalsize\color{black!80}
    \bulletitem[mygreen]{Chaînes de Markov et méthode Monte-Carlo}
    \bulletitem[mygreen]{Cas d'usage : détection de fissures. Petit historique : un retour sensible au \emph{model-based} ?}
    }
};
\adddatechap{C9}{1\textsuperscript{er} -- 10 mai}

\node[chapbox, below=0.3cm of C9.south west, anchor=north west] (C10) {
    \textbf{\large Chapitre 10 : Des (hyper)graphes plus riches (\emph{features}, données manquantes, \&c.)}\\
    \vspace{1.5mm}
    {\normalsize\color{black!80}
    \bulletitem[mygreen]{Un cadre bayésien à la détection de communautés sur graphes}
    \bulletitem[mygreen]{Les dynamiques de clusters d'ordre supérieur}
    \bulletitem[mygreen]{Un exemple où la percolation d'ordre supérieur affine des bornes théoriques}
    }
};
\adddatechap{C10}{10 -- 15 mai}

\node[chapbox, below=0.3cm of C10.south west, anchor=north west] (C11) {
    \textbf{\large Chapitre 11 : Élargissement au traitement du signal : le cas d'usage des fissures sur images}\\
    \vspace{1.5mm}
    {\normalsize\color{black!80}
    \bulletitem[mygreen]{Un classique : le filtre de Frangi}
    \bulletitem[mygreen]{Une structure plus riche : le graphe de Frangi}
    \bulletitem[mygreen]{Résultats}
    \bulletitem[mygreen]{Ouverture : vers une réconciliation de l'apprentissage profond et des modèles}
    }
    \vspace{1.5mm}
    {\small\color{myorange}\faQuestionCircle\ \textit{Question : Parler des travaux de Jules et de sa suite ?}}
};
\adddatechap{C11}{15 -- 20 mai}

\node[partbox=mypurple, below=0.8cm of C11.south west, anchor=north west] (IC) {
    \textbf{\Large \color{mypurple}Introduction générale \& Conclusion générale}\\
    \vspace{1.5mm}
    {\normalsize\color{black!80}
    \bulletitem[mypurple]{Propos liminaire}
    \bulletitem[mypurple]{Une brève histoire du \emph{clustering}. L'omniprésence du \emph{Single-Linkage}}
    \bulletitem[mypurple]{Les deux grandes familles de \emph{clustering} (besoin de vue hiérarchique)}
    }
    \vspace{1.5mm}
    {\small\color{myorange}\faEdit\ \textit{Action : Reprendre une partie très synthétique du rapport de recherche de fév. 2026.}}
};
\adddatepart{IC}{20 mai\\-- 1\textsuperscript{er} juin}{mypurple}

\node[milebox, below=1.2cm of IC.south west, anchor=north west] (M1) {
    \textbf{\Large \color{myorange}\faBook\ \ Livraison du Manuscrit aux Rapporteurs}
};
\adddatemile{M1}{15 juin}

\node[milebox, below=0.6cm of M1.south west, anchor=north west] (M2) {
    \textbf{\LARGE \color{myorange}\faGraduationCap\ \ Soutenance de Thèse !}
};
\adddatemile{M2}{8 sept.}

\begin{scope}[on background layer]
    \draw[gray!20, line width=6pt] (0, 1.5) -- (0, 0 |- M2.center);
\end{scope}
\end{tikzpicture}

\end{document}